\subsection{Probing Resources, Memory Resolution and Monetary Output}
\label{sec:resource_frontier}
The central model question is which resource buys identification. Observation length alone is insufficient when the lifetime quantity of controlled executions is fixed. This connects the merged-memory experiment to a financially specified input/output contract, rather than equating a state-norm shock with executed volume. Here $\Phi$ is the standard normal distribution function.

\begin{theorem}[Sharp probing-resource orders for a known symmetric pair]
\label{thm:resource_frontier}
Fix $0\le a<1$, $\kappa,q>0$, $S=a+\kappa q<1$, and $\eta\in(0,1)$. Compare a single cohort of weight $q$ and rate $\eta$ with two cohorts of weights $q/2$ and rates $\eta\pm\delta$. Require both split rates to lie in $[\epsilon,1-\epsilon]$, with $\epsilon>0$. All parameters of the two alternatives are known. The price recursion receives executed quantity $v_t$ through $\kappa v_t$, starts from zero, has no process noise, and is observed with independent $N(0,\sigma^2)$ measurement noise, $\sigma>0$. Probes are deterministic and common to both alternatives. Let $h_j=\kappa e_x^\top J^j(1,\eta_1,\ldots,\eta_K)^\top$ denote the impulse response.

Set $K_*=2\kappa^2q/[\epsilon^3(1-S)^2]$. For every finite horizon and every probe with total quantity $\sum_t|v_t|\le V$,
\begin{equation}
D_{\rm KL}(P_0^v\|P_\delta^v)
\le \frac{K_*^2 V^2\delta^4}{2\sigma^2},\qquad
e_{\rm worst}\ge\frac12\left(1-\frac{K_*V\delta^2}{2\sigma}\right)_+.
\label{eq:quantity_information}
\end{equation}
More precisely, with observations $x_1,\ldots,x_T$, define $d_T=(h_{\delta,j}-h_{0,j})_{j=0}^{T-1}$. The optimal worst-case error over these deterministic total-quantity probes and all classifiers is exactly
\begin{equation}
e^*_T(V)=\Phi\!\left(-\frac{V\|d_T\|_2}{2\sigma}\right).
\label{eq:exact_quantity_error}
\end{equation}
A pulse at the first input, followed by the Gaussian likelihood-ratio test using all observations, attains this value. Waiting adds impulse-response information but $\|d_T\|_2\uparrow\|d_\infty\|_2<\infty$; it cannot make the error arbitrarily small at fixed $V$.
Consequently error at most $\alpha<1/2$ requires
$V\ge2\sigma(1-2\alpha)/(K_*\delta^2)$, independently of observation length. Conversely, one pulse of magnitude
\begin{equation}
A=\frac{2\sigma\Phi^{-1}(1-\alpha)}{\kappa^2q\delta^2}
\label{eq:quantity_attainability}
\end{equation}
and the third price observation achieve worst-case error exactly $\alpha$. Thus the necessary and sufficient total-quantity order is $\delta^{-2}$ for this known pair.

Under average power $\sum_{t=0}^{T-1}v_t^2\le PT$, $P>0$, the same experiment instead has necessary horizon order $\delta^{-4}$, attained by the pulse above whenever $T\ge\max\{3,\lceil A^2/P\rceil\}$. This does not solve estimation of an unknown mixture.

Finally, for a fixed monetary price exposure $\lambda\ne0$, an additive reserve $L$ and cost envelope $c\sum|v_t|$, $c\ge0$, the sharp robust total-quantity allowance under budget $B>L$ is
\begin{equation}
V_B=\frac{B-L}{|\lambda|\kappa+c}.
\label{eq:financial_resolution}
\end{equation}
Within that allowance, error at most $\alpha$ requires
$\delta^2\ge2\sigma(1-2\alpha)/(K_*V_B)$; the pulse construction suffices if $A\le V_B$. These are distinct lower and upper constants, not an exact optimal-test threshold or a universal financial safety ceiling.
\end{theorem}

\begin{proof}
For either positive economy, $J$ is nonnegative and its price row sums to $S$. Cohort row $i$ sums to $1-\eta_i+\eta_i S$. Hence $\|J\|_\infty\le1-\epsilon(1-S)<1$, and the impulse response is nonnegative. Its maximum is $h_0=\kappa$, and its sum is $H(1)=\kappa/(1-S)$.

Write $G_0(z)=q\eta z/(z-1+\eta)$ and
$G_\delta(z)=(q/2)\sum_{s=\pm1}(\eta+s\delta)z/(z-1+\eta+s\delta)$.
Direct algebra gives
\[
G_\delta-G_0
=-\frac{qz(z-1)\delta^2}{(z-1+\eta)((z-1+\eta)^2-\delta^2)}.
\]
On $|z|=1$, each of the three denominator factors has modulus at least $\epsilon$, so $|G_\delta-G_0|\le2q\delta^2/\epsilon^3$. Since $H(z)=\kappa/[z-a-\kappa G(z)]$, the identity $H_\delta-H_0=H_\delta H_0(G_\delta-G_0)$ and positivity give
$\|H_\delta-H_0\|_\infty\le K_*\delta^2$.
Parseval bounds $\|h_\delta-h_0\|_2$ by the same quantity. Young's convolution inequality yields $\|\mu_\delta-\mu_0\|_2\le K_*\delta^2\|v\|_1$; Gaussian equal-covariance KL and Pinsker give \eqref{eq:quantity_information}. For average power, the induced $\ell_2$ gain instead gives $D_{\rm KL}\le K_*^2\delta^4PT/(2\sigma^2)$, proving the horizon lower order.

For the exact finite-horizon statement, the mean difference is a lower-triangular convolution matrix times $v$. Each column is a shorter suffix-shifted copy of $d_T$, with norm at most $\|d_T\|_2$. The triangle inequality therefore gives $\|\mu_\delta-\mu_0\|_2\le V\|d_T\|_2$, with equality for the first pulse. For two known equal-covariance Gaussian laws, their minimax error is $\Phi[-\|\mu_\delta-\mu_0\|_2/(2\sigma)]$, attained by the midpoint likelihood-ratio test. This proves \eqref{eq:exact_quantity_error} including the indistinguishable first-two-observation case.

The first two impulse taps are equal. Multiplying the price-first cohort transitions twice gives
$h_{\delta,2}-h_{0,2}=-\kappa^2q\delta^2$. A pulse $A$ therefore produces a third-observation mean separation $A\kappa^2q\delta^2$. The Gaussian midpoint test has error $\Phi[-A\kappa^2q\delta^2/(2\sigma)]$, proving attainability for both resource protocols. Extending the observation horizon does not create additional executed quantity.

Superposition and $\max_jh_j=\kappa$ imply $\sup_t|\lambda x_t|\le|\lambda|\kappa V$. A single pulse attains the bound at its first response, simultaneously attaining the declared cost envelope, so \eqref{eq:financial_resolution} is sharp for that robust output-plus-cost criterion. Substitution in the information bounds proves the final statements.
\end{proof}

The monetary reserve and cost functional are imposed model inputs, not estimated market losses. The result concerns an open-loop controlled experiment with no natural process excitation. A fixed lifetime quantity constraint, an average-power constraint and a repeatedly refreshed risk allowance are different experiments. In particular, without accumulated costs a maintained bounded-amplitude probe can continue; the theorem does not impose a capital ceiling on every safe adaptive strategy. Applying it to public markets requires actual executions, exposures, model-error coverage and a valid observation law.

The paper-specific advance proposed here is the explicit connection between memory resolution, probing resources and monetary output, with matching orders for a known pair. Partial fractions, Gaussian testing, convolution inequalities and positive-system gains are established tools. Closely spaced source recovery is already studied in mixture and super-resolution theory \citep{manoleho2022,liuammari2023}; non-normal financial bubbles and Kreiss diagnostics also precede this paper \citep{sornette2023nonnormal}. These comparisons exclude first-use claims. Independent assessment of the importance and novelty of this particular connection remains necessary.
